\noindent
At this point we start discussing the proper base change theorem.
To do so we introduce some notation. consider a commutative diagram
\begin{equation}
\label{equation-base-change-diagram}
\vcenter{
\xymatrix{
X' \ar[r]_{g'} \ar[d]_{f'} & X \ar[d]^f \\
Y' \ar[r]^g & Y
}
}
\end{equation}
of morphisms of schemes. Then we obtain a commutative diagram of sites
$$
\xymatrix{
X'_\etale \ar[r]_{g'_{small}} \ar[d]_{f'_{small}} &
X_\etale \ar[d]^{f_{small}} \\
Y'_\etale \ar[r]^{g_{small}} &
Y_\etale
}
$$
For any object $E$ of $D(X_\etale)$ we obtain a canonical base change map
\begin{equation}
\label{equation-base-change}
g_{small}^{-1}Rf_{small, *}E \longrightarrow Rf'_{small, *}(g'_{small})^{-1}E
\end{equation}
in $D(Y'_\etale)$. See Cohomology on Sites, Remark
\ref{sites-cohomology-remark-base-change} where we use the constant
sheaf $\mathbf{Z}$ as our sheaf of rings.
We will usually omit the subscripts ${}_{small}$ in this formula.
For example, if $E = \mathcal{F}[0]$ where $\mathcal{F}$ is an abelian
sheaf on $X_\etale$, the base change map is a map
\begin{equation}
\label{equation-base-change-sheaf}
g^{-1}Rf_*\mathcal{F} \longrightarrow Rf'_*(g')^{-1}\mathcal{F}
\end{equation}
in $D(Y'_\etale)$.

\medskip\noindent
The map (\ref{equation-base-change}) has no chance of being an isomorphism
in the generality given above. The goal is to show
it is an isomorphism if the diagram (\ref{equation-base-change-diagram})
is cartesian, $f : X \to Y$ proper, the cohomology sheaves of $E$
are torsion, and $E$ is bounded below.
To study this question we introduce the following terminology.
Let us say that {\it cohomology commutes with base change
for $f : X \to Y$} if (\ref{equation-base-change-sheaf})
is an isomorphism for every diagram (\ref{equation-base-change-diagram})
where $X' = Y' \times_Y X$ and every torsion abelian sheaf $\mathcal{F}$.
